%======================================================================
\section{The full language}\label{sec:full}

In this section formulas are built from variables and $\bot$ by $\wedge$,
$\vee$ and $\to$. For a nonempty list $\Delta=(\delta_1,\dots,\delta_k)$ we
write $\bigwedge\Delta$ for $\delta_1\wedge(\delta_2\wedge\dots\wedge\delta_k)$.
\kl{Elementary generalizations} now include \kl{abstractions} of sets of
\kl{occurrences} of $\bot$, and Proposition~\ref{prop:criterion} applies as
stated.

A formula of $\CPC$ with a \kl{positive} \kl{occurrence} of $\bot$ is never
\kl{minimal in} $\CPC$: replacing that \kl{occurrence} by a new variable $z$
gives a formula implied by the original one in $\IPC$, by ex falso and
monotonicity, and this formula is an \kl{abstraction}. Our construction
therefore keeps $\bot$ only in \kl{negative} positions.

\subsection{Local constraints and the monitor}

Let $\Phi\in\CPC$; then $\Phi$ is neither a variable nor $\bot$. The
\kl{local constraints} are defined as in Section~\ref{sec:constraints}:
variables are \kl{renamed apart} with their \kl{guards}, every
\kl{non-leaf occurrence} receives an \kl{extension variable}, and so does
every \kl{occurrence} of $\bot$. Writing $\nv{v}$ for the \kl{node variable}
of $v$, with $u=v0$ and $w=v1$, the constraint $\cons{v}$ is
\[
\begin{array}{lll}
\subf{\Phi}{v} & \text{$v$ \kl{positive}} & \text{$v$ \kl{negative}}\\\hline
\text{a variable }p & p\to\nv{v} & \nv{v}\to p\\
\subf{\Phi}{u}\circ\subf{\Phi}{w},\ \circ\in\{\wedge,\vee,\to\}
 & (\nv{u}\circ\nv{w})\to\nv{v} & \nv{v}\to(\nv{u}\circ\nv{w})\\
\bot & \text{(none)} & \nv{v}\to\bot
\end{array}
\]
No constraint is needed for a \kl{positive} $\bot$, since $\bot\to\nv{v}$ is
provable. Let $W$ be the set of \kl{occurrences} that have a constraint, and
$W_\bot\subseteq W$ the set of \kl{negative} \kl{occurrences} of $\bot$.
Every \kl{local constraint} has \kl{order} at most $3$, and every one not of
the form $\nv{v}\to\bot$ is \kl{read-once} without $\bot$. \AP For a
valuation $\eta$ of the variables of $\Phi$, the \intro{standard valuation}
$\intro*\stdval{\eta}$ extends $\eta$ by
$\stdval{\eta}(\nv{v})=\eta(\subf{\Phi}{v})$.

\begin{lemma}\label{lem:full-constraints}
The \kl{local constraints} derive $\Phi\to\nv{\rt}$ in $\IPC$. The
substitution $\nv{v}\mapsto\subf{\Phi}{v}$, fixing the variables of $\Phi$,
maps every \kl{local constraint} to a formula $\gamma\to\gamma$;
consequently every \kl{local constraint} is true under every
\kl{standard valuation}.
\end{lemma}
\begin{proof}
As in Lemmas~\ref{lem:port} and~\ref{lem:S-sound}, using the monotonicity of
$\wedge$ and $\vee$; a \kl{positive} $\bot$ is handled by $\bot\to\nv{v}$.
\end{proof}

The \kl{monitor} is modified so as not to copy $\bot$ into a \kl{positive}
position. We take \kl{selectors} $\sel{v}$ $(v\in W)$, a \kl{switch} $\sw$,
and new variables $\intro*\fresh{v}$ $(v\in W_\bot)$, and put
$\guarded{v}=\sel{v}\to\cons{v}$, $\link{v}=\sw\to\sel{v}$,
$\query=\sw\to\nv{\rt}$ as before, and
\[
\widehat{\cons{v}}=\nv{v}\to\fresh{v}\ (v\in W_\bot),\qquad
\widehat{\cons{v}}=\cons{v}\ (v\notin W_\bot),\qquad
\widehat{\guarded{v}}=\sel{v}\to\widehat{\cons{v}} .
\]
With $\Prem$ the list of the $\guarded{v},\link{v}$ and $\widehat{\Prem}$
the list of the $\widehat{\guarded{v}},\link{v}$ $(v\in W)$, let
\[
\Fst{\Phi}=\bigwedge\Prem\to\Bigl(\query\wedge\bigwedge\widehat{\Prem}\Bigr).
\]
The conjunction replaces the continuation $(\query\to\Prem\to\goal)\to\goal$
of Section~\ref{sec:localization}, which lowers the \kl{order} by one. The
\kl{left} and \kl{right} members are those in $\Prem$ and in
$\widehat{\Prem}$. Each $\fresh{v}$ occurs exactly once, and $\bot$ occurs
only in \kl{left} members, \kl{negatively}.

\begin{lemma}\label{lem:full-F}
\leavevmode
\begin{enumerate}
\item[(a)] $\IPC\vdash\Phi\to\Fst{\Phi}$, and
  $\IPC\vdash\elimsub{\Phi}(\Fst{\Phi})\to\Phi$ for the substitution
  $\intro*\elimsub{\Phi}$ that fixes the variables of $\Phi$ and maps
  $\nv{v}\mapsto\subf{\Phi}{v}$, $\sw,\sel{v}\mapsto\ttop$ and
  $\fresh{v}\mapsto\bot$.
\item[(b)] Every variable of $\Fst{\Phi}$ other than the $\fresh{v}$ has a
  \kl{positive} and a \kl{negative} \kl{occurrence}.
\item[(c)] No \kl{separation} of $\Fst{\Phi}$ belongs to $\CPC$.
\item[(d)] No \kl{abstraction} of a nonempty set of \kl{occurrences} of
  $\bot$ in $\Fst{\Phi}$ belongs to $\CPC$.
\end{enumerate}
\end{lemma}
\begin{proof}
(a) As in Lemma~\ref{lem:F-equiv}, using Lemma~\ref{lem:full-constraints};
the \kl{right} member $\sel{v}\to\nv{v}\to\fresh{v}$ follows from
$\sel{v}\to\nv{v}\to\bot$ by ex falso. Under $\elimsub{\Phi}$ the members of
$\Prem$ become theorems, and the first conjunct of the consequent becomes
$\ttop\to\Phi$.

(b) As in Lemma~\ref{lem:F-polar}; a \kl{node variable} of a \kl{positive}
$\bot$ occurs in the constraint of its parent.

(c) We follow the proof of Lemma~\ref{lem:F-sep}, with every $\fresh{v}=0$,
so that $\widehat{\cons{v}}$ and $\cons{v}$ have the same value. Now
$\beta$ is false as soon as all \kl{left} members are true and $\query$ or
some \kl{right} member is false, and case~(h) disappears. In cases
(a)--(d) we use Lemma~\ref{lem:readonce}, and for $\cons{v}=\nv{v}\to\bot$
the fact that its value is the complement of that of $\nv{v}$. In cases
(e) and~(g), instead of giving all \kl{data variables} the value $1$, we give
them the values of a \kl{standard valuation}, under which every \kl{local
constraint} is true by Lemma~\ref{lem:full-constraints}.

(d) Let the chosen set contain the $\bot$ of $\guarded{v}^L=\sel{v}\to\nv{v}\to\bot$
and let it be replaced by $z$. Let $z=1$, $\sw=0$, $\sel{v}=\nv{v}=1$,
$\sel{w}=0$ for $w\ne v$, and $\fresh{v}=0$. All \kl{left} members are true,
whereas $\widehat{\guarded{v}}=\sel{v}\to\nv{v}\to\fresh{v}$ is false.
\end{proof}

\subsection{The tree closure}

Let $F$ be a formula in which every variable that occurs at least twice has
a \kl{positive} and a \kl{negative} \kl{occurrence}. We \kl{rename apart}
the \kl{occurrences} of these variables only, leaving $\bot$ and the
variables that occur once unchanged in $\ported{F}$, and define
\kl{guards} and \kl{edge formulas} as in Section~\ref{sec:minimality}. In
this section the \kl{tree closure} of $F$ is
\[
\Tc{F}=\bigwedge\Guards_F\to\Bigl(\ported{F}\wedge\bigwedge\Edges_F\Bigr).
\]

\begin{theorem}\label{thm:full-closure}
Let $F\in\CPC$ be as above, and suppose that no \kl{separation} of $F$ and
no \kl{abstraction} of a nonempty set of \kl{occurrences} of $\bot$ in $F$
belongs to $\CPC$, and that every subformula $a\circ b$ of $F$ with $a,b$
variables and $\circ\in\{\wedge,\vee,\to\}$ contains a variable that occurs
at least twice in $F$. Then $\Tc{F}$ is \kl{minimal in} $\CPC$,
$\IPC\lplus\Tc{F}=\IPC\lplus F$, and $\Ord(\Tc{F})\le\max\{3,\Ord(F)\}$.
\end{theorem}
\begin{proof}
The equivalence is proved as in Lemma~\ref{lem:M-equiv}, using the
monotonicity of $\wedge$ and $\vee$.

\kl{Splittings}. A \kl{splitting} replaces \kl{occurrences} of a
\kl{center} or of a \kl{port}. With ``$\ported{F}$ or some \kl{edge formula}
is false'' as the condition for falsity, cases (a)--(c) of the proof of
Lemma~\ref{lem:M-split} apply verbatim; the values of the variables that
occur once are taken from the \kl{separation} in case~(c) and are arbitrary
otherwise.

\kl{Abstractions} of $\bot$. The \kl{occurrences} of $\bot$ in $\Tc{F}$
are those in $\ported{F}$. Given an \kl{abstraction}, take a falsifying
valuation of the corresponding \kl{abstraction} of $F$, give every
\kl{port} the value of its \kl{center}, and keep the values of the other
variables. All \kl{guards} and \kl{edge formulas} are true, and the
modified $\ported{F}$ is false.

\kl{Abstractions} of compound formulas. The proof of
Lemma~\ref{lem:abs-necessary} applies, with one addition: if the
abstracted formula $\gamma$ contained $\bot$ at a \kl{leaf} $\ell$, then
substituting for $z$ the formula $\gamma$ with that $\bot$ replaced by a
new variable would give an \kl{abstraction} of $\bot$ in $\CPC$. Hence there
are distinct variables $a,b$ and $\circ\in\{\wedge,\vee,\to\}$ such that
all \kl{occurrences} of $a$ and $b$ lie in \kl{occurrences} of $a\circ b$.
As in Lemma~\ref{lem:M-noabs}, this is impossible: such a subformula is a
\kl{guard}, or an \kl{edge formula}, or lies in $\ported{F}$ and then
contains a \kl{port} by the last hypothesis; and every
\kl{port} also occurs in its \kl{guard}. The conjunctions joining the
members, and the main implication, connect compound formulas.

By Proposition~\ref{prop:criterion}, $\Tc{F}$ is \kl{minimal in} $\CPC$.
\kl{Guards} and \kl{edge formulas} have \kl{order} $2$, so
$\Ord(\Tc{F})\le\max\{3,\Ord(F)\}$.
\end{proof}

Let $\Mz{\Phi}=\Tc{\Fst{\Phi}}$.

\begin{theorem}\label{thm:full}
For every $\Phi\in\CPC$, the formula $\Mz{\Phi}$ is \kl{minimal in} $\CPC$,
$\Ord(\Mz{\Phi})\le4$, $\IPC\vdash\Phi\to\Mz{\Phi}$ and
$\IPC\vdash\recsub{\Phi}(\Mz{\Phi})\to\Phi$ for a substitution
$\recsub{\Phi}$ fixing the variables of $\Phi$; hence
$\IPC\lplus\Mz{\Phi}=\IPC\lplus\Phi$. Moreover, if $\Phi$ has $B_+$
\kl{positive} and $B_-$ \kl{negative} \kl{occurrences} of $\bot$, then
$\sizeof{\Mz{\Phi}}=102\sizeof{\Phi}-48B_+-8B_--2|\Var(\Phi)|-48$.
\end{theorem}
\begin{proof}
By Lemma~\ref{lem:full-F} and Theorem~\ref{thm:full-closure}, with
$\recsub{\Phi}=\elimsub{\Phi}\circ\collsub$; the last hypothesis of
Theorem~\ref{thm:full-closure} holds, since the only variables occurring once
in $\Fst{\Phi}$ are the $\fresh{v}$, which occur in
$\nv{v}\to\fresh{v}$. Every \kl{local constraint} has
\kl{order} at most $3$, so $\bigwedge\Prem$ has \kl{order} at most $3$ and
$\Ord(\Fst{\Phi})\le4$. For the size, let $L=\sizeof{\Phi}$ and
$V=|\Var(\Phi)|$. There are $n=2L-1-B_+$ \kl{local constraints}, of total
size $5L-2B_+-3$, so $f:=\sizeof{\Fst{\Phi}}=2(11L-5B_+-6)+2$. The number of
\kl{ports} is $m=f-2B_-$ and the number of \kl{centers} is
$w=V+(2L-1)+n+1$. Hence $\sizeof{\Mz{\Phi}}=2m+f+2(m-w)$, which gives the
formula.
\end{proof}

\begin{proof}[Proof of Theorem~\ref{thm:intro-full}]
Take $M=\Mz{\Phi}$ of Theorem~\ref{thm:full}. Every intermediate logic $L$
equals $\IPC\lplus L$, and replacing each axiom $\Phi\in L$ by $\Mz{\Phi}$
gives an axiomatization by formulas \kl{minimal in} $\CPC$.
\end{proof}

\plot{Hilbert-style axiomatization of $\IPC$ by formulas minimal in $\CPC$
(check each axiom), for the analogue of Remark~\ref{rem:base}.}

\subsection{Lower orders}

In the full language the bound $4$ is not optimal for every logic.

\begin{example}\label{ex:order2}
The formula $p\vee(p\to q)$ has \kl{order} $2$, is \kl{minimal in} $\CPC$,
and axiomatizes $\CPC$ over $\IPC$: substituting $\bot$ for $q$ gives
$p\vee\neg p$, and conversely $p\vee\neg p$ implies $p\vee(p\to q)$ by ex
falso. Its \kl{elementary generalizations} are refuted as follows: the two
\kl{splittings} of $p$ by giving the first \kl{occurrence} of $p$ the value
$0$, the second the value $1$, and $q=0$; the \kl{abstraction}
$p\vee z$ by $p=z=0$; and $z$ by $z=0$.
\end{example}

\begin{example}\label{ex:order3}
The formula $(p\to q)\vee(\neg p\to r)$ has \kl{order} $3$, is \kl{minimal
in} $\CPC$, and axiomatizes Jankov's logic $\intro*\KC=\IPC\lplus(\neg p\vee\neg\neg p)$~\cite{Jan68b}:
substituting $\bot$ for $q$ and $r$ gives $\neg p\vee\neg\neg p$, and
conversely each disjunct of $\neg p\vee\neg\neg p$ implies a disjunct of
the formula by ex falso. The \kl{elementary generalizations} are refuted as
follows: the two \kl{splittings} of $p$ by giving the first \kl{occurrence}
of $p$ the value $1$, the second the value $0$, and $q=r=0$; abstracting $\bot$ by $p=z=1,q=r=0$; abstracting
$p\to q$ by $z=p=r=0$; abstracting $\neg p$ by $p=z=1,q=r=0$; abstracting
$\neg p\to r$ by $p=1,q=z=0$; and abstracting the whole formula.
\end{example}

\begin{remark}\label{rem:bot}
$\bot$ cannot be eliminated: $\KC$ and $\IPC$ have the same $\bot$-free
theorems~\cite{Jan68b}, so no $\bot$-free formula $\Psi$ satisfies
$\IPC\lplus\Psi=\KC$. In particular, Theorem~\ref{thm:intro-main} does not
yield Theorem~\ref{thm:intro-full}.
\end{remark}
